\documentclass[12pt, a4paper]{article}
\usepackage{graphicx} % Required for inserting images
\usepackage{amsmath}
\usepackage{amsfonts}
\usepackage{amssymb}
%\usepackage[ruled, vlined]{algorithm2e}
\usepackage{algorithm}
\usepackage{algorithmic}
\usepackage{enumitem}

\textwidth=150mm
\textheight=250mm
\hoffset= -10mm       % may need change
\voffset= -25mm 

\title{Finding the ideal factor for adding the defect correction in a multigrid solver for the Laplician equation with discontinuous coefficients}
\author{Jakob Niessner}
\date{\today}

\begin{document}

\maketitle

\newpage

\tableofcontents

\newpage

\section{Theoretical background}

The multigrid method in its most general form is outlined below. We denote $u_i,$ the current guess, $A^{h/H}$  the differential operator on the fine ($h$) or on the coarse ($H$) grid, and $f$  the right hand side.

\begin{algorithm}[H]


\caption{Multigrid cycle}

\vspace{0.4cm}
\begin{itemize}
    \item Pre-smoothing to reduce high-frequency errors.
     \begin{equation}
       v_i = S_{\textit{pre}}(u_i, f, A, \nu_1),
     \end{equation}
    \item Compute the residual on the fine grid:
        \begin{equation}
          r^h = f - A^h v_i,
        \end{equation}
        \item Restrict the residual according to some Restriction operator $R$
        \begin{equation}
          r^H = R r^h,
        \end{equation}
        \item Solve on the coarser grid (Possible using nestest multigrid cycles)
        \begin{equation}
          A^H e^H = r^H,
        \end{equation}
        \item Prolongate and correct. 
          \begin{equation}
            \tilde{v}_i = v_i + P e^H,
         \end{equation}
        \item Post-smoothing:
          \begin{equation}
            u_{i+1} = S_{\textit{post}}(\tilde{v}_i, f, A, \nu_2),
          \end{equation}
    \item Repeat until convergence criteria are met (e.g., residual norm below a threshold).
\end{itemize}

\end{algorithm}

\noindent
This scheme with a 2-iteration Gauss-Seidel smoother work very nicely if $A$ is a discretization of the Laplacian operator $\Delta = -\nabla \cdot \nabla$ on a uniform grid (Q1, Finite difference, etc. ). However, this can drastically change if $A$ is represents the Laplacian with discontinuous coefficients, i.e., $-\nabla \cdot (\epsilon \nabla)$ with $\epsilon$ being piecewise constant. In this case, the above scheme converges much worse up to the point were it diverges. It has therefore been proposed to add a defect correction term in step 5, i.e.,

\begin{equation}
  \tilde{v}_i = v_i + \omega P e^H 
\end{equation}


\noindent
The remaining question is: What is the best choice $\omega_{\textit{ideal}}$? The easiest starting point would be to simply guess $\omega_{\textit{ideal}}.$ However, this this is very tedious and probably very problem dependent. It is therefore desirable to find a more systematic way to determine it.  

\newpage

\subsection{Formula for the optimal $\omega$}

In order to find a formula for the optimal $\omega,$ we can start with a geometic argument. Note that if $\epsilon$ is positive, which it is, for example, in electrostatics. The operator $A$ is positive definite and symmetric. Therefore, the problem: Fin $x$ such that 

\begin{equation}
  Ax = f
\end{equation}


\noindent
is equivalent to the the optimization problem: Find $x$ such that the functional
\begin{equation}
  \Phi(x) = \frac{1}{2} \langle x,  A x \rangle - \langle x , f \rangle
\end{equation}
is minimial. Here, $\langle \cdot, \cdot \rangle$ denotes the standard inner product. Given that $\tilde{v}_i$ is the next guess before the post-smoothing, $\omega$ should be chosen such that $\Phi(\tilde{v}_i)$ is minimal. Hence, we should choose $\omega$ such that the function (denoting $e = e^H$ for simplicity)

\begin{equation}
 \phi(\omega) = \frac{1}{2} \langle v_i + \omega e,  A (v_i + \omega  e) \rangle - \langle v_i + \omega  e , f \rangle,
\end{equation}

\noindent
is minimal. Since the above function is a parabola, with a positive prefactor, the minimum is unique and given by the solution of the equation 

\begin{equation}
  \frac{d \phi}{d \omega} = 0.
\end{equation}

\noindent
This solution to this equaiton is given by
\begin{equation}
  \omega_{\textit{ideal}} = \frac{\langle f - A v_i, e \rangle}{\langle A e, e \rangle}.
  \label{eq:omega_ideal}
\end{equation}


\newpage

\section{Numerical experiments}

In order to find the optimal $\omega,$ we followed $2$ approaches. 
\begin{itemize}
  \item[1.)] First We guessed and tried to systematically narrow down to the optimal $\omega.$
  \item[2.)] We used the equation \eqref{eq:omega_ideal} to compute the ideal $\omega$.
\end{itemize}

\noindent
As a test case, we used a 3D-cube as a domain with a side length of $16$ units with a sphere of radius $1.8$ units in the center. The coefficient $\epsilon$ was set to $4$ inside the sphere and $8$ outside. As a FEM-discretization, we used a cubic mesh, where each cube was divided into $6$ tetrahedra, such that the resulting stencil is a $7$-point stencil. We used $513 \times 513 \times 513$ grid points, s.t. we hade $511$ degrees of freedom in eacht direction. The right hand side was set to $0$ everywhere, and the boundary conditions were set according to the function $u(x, y, y) = -3x^2 - 4y^2+7z^2$. $2$-pre and post-smoothing steps were used. And 2 levels, such that the coarse grid had $255 \times 255 \times 255$ degrees of freedom.


\subsection{Systematic search for the optimal $\omega$}

\subsection{Using the formula for $\omega_{\textit{ideal}}$}





\end{document}
